\section{Introduction}\label{sec:intro}

Photorealistic and animatable digital humans are a cornerstone of modern
multimedia, powering virtual production, immersive telepresence, augmented
reality try-on, and interactive gaming. A long-standing goal in this domain
is to create such avatars from the most ubiquitous capture modality---a
single monocular RGB video---without requiring a calibrated multi-view
studio. Recent progress on 3D Gaussian Splatting (3DGS)~\cite{kerbl3Dgaussians}
and on feed-forward 3D generation has brought this goal within reach, yet a
satisfactory solution that is simultaneously \emph{fast}, \emph{drivable},
and \emph{incrementally refinable} remains elusive.

Existing methods cluster into two paradigms with complementary strengths.
\textbf{Feed-forward generators}~\cite{yan2026triposplat,tang2024lgm,zhang2024gslrm,pan2024humansplat,qiu2025lhm}
map one or a few posed images to a set of 3D Gaussians in a single forward
pass, producing a usable asset in seconds. Their accuracy on unseen regions
and fine appearance detail is, however, fundamentally bounded by the
generalization capacity of the generator. \textbf{Per-video optimizers}
based on deformable or articulated 3DGS~\cite{hu2024gauhuman,hu2024gaussianavatar,qian2024gsavatar,kocabas2024hugs}
fit an avatar to an entire monocular sequence through gradient descent,
attaining high fidelity but requiring tens of minutes of offline training
and, crucially, treating the video as a closed dataset: once optimized, the
model cannot ingest additional footage without retraining from scratch.
The recently emerging line of streaming 3DGS~\cite{wu2025streamsplat,li2025streamgs,girish2024queen}
reconstructs general scenes online, but does not target the articulated,
skeleton-drivable human setting nor explicitly address the catastrophic
forgetting that arises when an avatar optimized on one set of viewpoints is
later exposed to new ones.

This gap motivates a key question: \emph{can we obtain a drivable avatar
within seconds of the first frame and then continuously improve its detail
as the monocular video streams in, without forgetting what was already
learned?} Naively sequential fine-tuning on streaming clips fails at this
goal: new observations overwrite earlier ones, and the avatar degrades on
previously well-reconstructed viewpoints---the well-known catastrophic
forgetting phenomenon, now manifest in the geometry and appearance of 3D
Gaussians rather than classifier weights.

We address this problem by casting monocular avatar reconstruction as a
\textbf{continual-learning} problem over 3DGS parameters. Our framework,
\method{}, makes the following contributions:

\begin{itemize}
    \item \textbf{A unified warm-start \& continual-refinement pipeline.}
          \method{} bootstraps a drivable 3DGS avatar from a single frame
          via a feed-forward Gaussian generator and an SMPL-based skeleton
          registration, then refines it continuously as monocular clips
          arrive. To our knowledge this is the first framework to
          explicitly connect continual learning with drivable 3DGS human
          avatars.
    \item \textbf{An EWC prior for 3D Gaussian parameters.} We estimate a
          per-Gaussian Fisher information importance map and regularize
          updates of important parameters when learning from new clips,
          preserving previously acquired detail without freezing the model.
    \item \textbf{A viewpoint-stratified replay buffer} that stores a
          compact, pose-diverse set of historical frames and re-renders
          them jointly with incoming data, providing episodic-memory
          replay adapted to the articulated structure of the human body.
    \item \textbf{Pose-conditioned selective reactivation} that focuses
          densification and optimization budget on body regions that are
          under-observed in the current pose, improving sample efficiency
          on streaming data.
\end{itemize}

Experiments on monocular human video benchmarks show that \method{}
approaches the fidelity of offline optimizers while consuming the video
online, and that its continual-learning components reduce forgetting of
early viewpoints by a large margin compared to naive sequential
fine-tuning. We release code to support reproducible research.
